\section{Related literature}
\citet{lucca2015} report large average excess equity returns before scheduled FOMC announcements: about 49~bps over the 24 hours ending roughly 15 minutes before the announcement, across 131 scheduled meetings from September 1994 to March 2011 (as stated in the New York Fed staff report version). They find similar pre-announcement returns in other major equity indices but not in Treasuries or money market futures, and no comparable returns before other macroeconomic announcements. The effect is larger when the yield curve is flat, when the VIX is high and when past pre-FOMC returns were high.

\citet{hu2022} study overnight returns before nonfarm payrolls, ISM and GDP releases and propose that heightened uncertainty about the impact of upcoming news, proxied by the VIX, is priced and then resolves at the announcement. \citet{martello2018}, a 2018 master's dissertation at PUC-Rio (SSRN 3286745), studies state dependence of the FOMC premium on implied volatility: its introduction reports an average drift of 109~bps when implied volatility is above its prior median and 9.7~bps below, in a baseline sample of 60 scheduled meetings from September 2003 to March 2011. This is the closest antecedent to our state split, which is why we fix the threshold from the development sample rather than from the holdout.

A separate strand asks whether the drift reflects information leakage or monetary policy communication. \citet{cieslak2019} show that equity returns since 1994 concentrate in even weeks of the FOMC cycle and link this to informal communication between Fed officials and market participants. \citet{bernile2016} find abnormal E-mini S\&P~500 futures trading during FOMC news embargoes that predicts later policy surprises. \citet{kurov2019} find signs of informed trading before about nine of twenty market-moving announcements. If such channels operate, measured premia depend on communication and compliance regimes, which changed with the move to a press conference after every meeting from 2019 \citep{boguth2019}. That paper reports that announcements with a press conference move prices more and that the market risk premium is higher on press conference days, which matters here because every 2020--2026 meeting had one.

\citet{mclean2016} study 97 cross-sectional predictors and report portfolio returns 26\% lower out of sample and 58\% lower post-publication. They study cross-sectional predictors, not announcement-window returns, so we use them only as a benchmark for expected decay. Our holdout tests the state-dependent prediction on data that none of these papers used.